\head{Grades K--1: untangle, then decide}
Accept any drawing with no crossings and the same links; below is one answer each, and for Problems 3 and 5 a drawing with the fewest crossings (circled). Ask ``How do you know for sure?'' only when a child is ready.
\begin{center}
\fig{key_k_bowtie}{P1, first}\hfill\fig{key_k_kite}{P1, second}\hfill\fig{key_k_k4}{P1, third}\hfill\fig{key_k_hw23}{P2}\hfill\fig{key_k_hw33}{P3: 1 crossing}\\[6pt]
\fig{key_k_star}{P4, first}\qquad\qquad\fig{key_k_fan}{P4, second}\qquad\qquad\fig{key_k_prism}{P4, third}\qquad\qquad\fig{key_k_k5}{P5: 1 crossing}
\end{center}

\sub{Problem 1, page 1: three four-dot networks}
All three untangle. Each picture has one crossing, and moving one dot fixes it. The ring A--B--C--D becomes a square; in the second, A and D go on opposite sides of link B--C; in the third, where every pair is joined, any one dot goes inside the triangle of the other three. Hint: ``Which two links cross? What if one of their dots moves?'' Extension for the third: ``Can you add one more link?'' (No; every pair is already joined.)

\sub{Problems 2 and 3, page 2: houses and wells}
Two houses and three wells: yes, with one house above the row of wells and the other below it. Hint: ``What if one house moves to the other side of the wells?'' Three houses and three wells: no; one crossing is the fewest (key, P3). Expect many children to get down to one crossing and stop. ``We couldn't, and one crossing is the fewest we found'' is the expected K--1 answer; it is evidence, not a proof, and the 4--5 pages prove it. Hint: ``How few crossings can you get?'' Extension: ``Take away any one pipe cleaner. Can you untangle it now?'' (Yes, whichever one goes.)

\sub{Problem 4, page 3: five and six dots}
All three untangle; the pictures have 5, 3 and 4 crossings. The first is one ring, A--B--C--D--E. In the second, E is joined to every other dot and A--B--C--D is a path, so E goes above a row A, B, C, D. The third is two triangles, A, B, C and D, E, F, joined by A--D, B--E and C--F; one triangle goes inside the other. Hints: ``Start at A and follow the links. Which dot comes next?'' ``Can you find three dots all joined to each other?''

\sub{Problem 5, page 4: five dots, all joined}
One crossing is the fewest (key, P5); expect children to reach two or three first. Ten links is a lot of twisting: an adult may attach them while the child decides where the dots go. An argument an adult can show with the pipe cleaners: untangle A, B, C, D first, as in Problem 1's third network. Every region then has three dots round it, so wherever E goes it reaches only three of the four without a crossing. Hint: ``Can you untangle four of the dots first? Where can the fifth go?''

\sub{Problem 6, page 5: the link game}
With 4 dots the second player wins every game, and with 5 dots the first player does, however anyone plays: a finished game has 6 links on 4 dots and 9 on 5. On 4 dots it joins every pair; on 5 dots exactly one pair stays unjoined (compare Problem 5). The risk is stopping early, when a link between the inner dots or along the outer triangle still fits. Before a game ends, ask: ``Is there a pair of dots that is not joined and still can be?'' If everyone draws straight links, the game is over exactly when no straight link fits; once someone draws a curved link, check curved routes too, including round the outside. Either way the total is the same. Let each child play first in some games and second in others, so the tally shows that the winning seat, not the player, repeats. Extension: ``Count the links at the end of each game.''

\head{Grades 2--3: sort, count, and the portal board}
\begin{center}
\fig{key_n1}{Network 1: 5, 7, 4}\hfill\fig{key_n3}{Network 3: 5, 9, 6}\hfill\fig{key_n4}{Network 4: 6, 9, 5}\hfill\fig{key_n6}{Network 6: 6, 12, 8}\\[4pt]
{\small\textbf{Figure 1.} The networks that untangle, with dots, links and regions.}\\[10pt]
\fig{few_K5}{Network 2: 1 crossing}\qquad\fig{few_K33}{Network 5: 1 crossing}\qquad\fig{few_K6}{6 dots: 3 crossings}\\[4pt]
{\small\textbf{Figure 2.} Fewest crossings, circled. Network 2's drawing is network 3's plus D--E.}
\end{center}

\sub{Problem 1, pages 1--2: untangled or stuck}
Networks 1, 3, 4 and 6 untangle (Figure 1); networks 2 and 5 stay stuck. The pictures have 2, 3, 3, 4, 5 and 4 crossings. The pairs are deliberate: network 3 is network 2 without D--E; networks 4 and 5 both have 6 dots and 9 links; network 6 has the most links and still untangles. Network 5 is the houses and wells. At this point ``stuck'' is a conjecture. Problem 6 asks why for network 2; network 5 stays a conjecture in this packet, though a ready child can try the overview's ring argument. Hints: for network 6, ``Which dot is A not joined to?'' (F; B--D and C--E are also missing) and ``What if A, B and C made a big triangle?'' For network 3, ``Which pair of dots is not joined?''

\sub{Problem 2, page 3: dots, links and regions}
Dots $-$ links $+$ regions $=2$ for every untangled network in one piece; children may write ``regions $=$ links $-$ dots $+2$'' or ``dots $+$ regions $=$ links $+2$''. The rows are in Figure 1, and the example has 4, 5, 3. Likely slips: forgetting the outside, counting a picture that still has crossings, and losing track in a crowded drawing; shading each region as it is counted helps. Hint: ``What happens with just one link? With two?'' Extensions: test a network in two pieces (it gives 3), and count the regions of a tangled picture (each crossing adds one). For a child who asks why: add the links one at a time and watch the three numbers.

\sub{Problem 3, page 4: as few crossings as you can}
Network 2: 1. Network 5: 1. Six dots all joined: 3 (Figure 2). Untangle network 3 and add D--E, which crosses one link; untangle houses and wells with one link removed and put it back across one link. For six dots, try one triangle inside another. That 1 is the best for network 2 is Problem 6's question. That 1 is the best for network 5 and 3 for six dots stay conjectures here (4--5 Problems 5 and 7 prove them). Hint: ``Which link could you take away to untangle it?''

\sub{Problem 4, page 5: the link game}
4 dots: second player; 5 dots: first player; 6 dots: second player, in every game. Finished games have 6, 9 and 12 links. Check endings as in K--1 Problem 6: ``Is there a pair that is not joined and still can be?'' Extensions: ``Count the links at the end.'' ``Predict 7 dots, then test.'' (15 links; the first player wins.)

\begin{center}
\fig{portal_k33}{Network 5}\qquad\qquad\fig{portal_k5}{Network 2}\\[4pt]
{\small\textbf{Figure 3.} One portal-board drawing at the printed positions. As in the page's example, a link that leaves at a number comes back in at the same number on the opposite side.}
\end{center}

\sub{Problem 5, page 6: the portal board}
Both can be drawn (Figure 3). Network 5: A joins all three wells directly, as do B--F and C--F; B--D, B--E and C--E leave through the top and come back in at the bottom; C--D leaves through the right side and comes back on the left. Network 2: join each pair by its shortest route, counting routes through the sides; five links use a portal. Many other drawings work. Hints: ``Which link is blocked inside the board? Could it leave through a side?'' ``Where does it come back in?'' Extension: count regions on the portal board, following each region across the sides. Every crossing-free drawing gives 3 for network 5 and 5 for network 2, so dots $-$ links $+$ regions $=0$ here.

\sub{Problem 6, page 7: why network 2 is stuck}
Three explanations children can reach. (1) A, B, C and D are all joined. Untangled, every region has just three of these dots round it (one dot sits inside the triangle of the other three; in general $6=3\cdot4-6$, so every region is a triangle), and E, in whichever region it sits, reaches only those three. (2) Draw the ring A--B--C--D--E first; each of the other five links goes inside or outside it, two fit inside and two outside, and the fifth must cross. (3) Untangled, Problem 2's rule gives 7 regions, each needing at least 3 links round it: 21 link-sides, but 10 links have 20. ``We tried lots of ways'' is evidence; ask ``Could someone else find a way? How do you know?'' Accept a partial argument with a drawing.

\head{Grades 4--5: proofs, the game and the portal board}
\sub{Problems 1 and 2, pages 1--3}
As Grades 2--3 Problems 1 and 2, without the split: the pair does all six networks and may draw instead of building.

\sub{Problem 3, page 4: the link game}
As Grades 2--3 Problem 4: second, first, second player. Collect the tally now; Problem 8 asks why.

\sub{Problems 4 and 5, page 5: why networks 2 and 5 are stuck}
Network 2: if it were untangled, Problem 2's rule would give $10-5+2=7$ regions. With no double links, each region has at least 3 links round it, so at least 21 link-sides, but 10 links have only 20. The ring and four-dot arguments of Grades 2--3 Problem 6 are also complete proofs. Follow-up: ``Why at least three links round a region?'' (Two would mean a double link.)
Network 5: every link joins a house to a well, so a loop alternates houses and wells and needs at least 4 links; no three dots are all joined. Untangled: $9-6+2=5$ regions, at least 4 link-sides each, so 20, but 9 links have 18. Ring version: around A--D--B--E--C--F the links A--E, B--F and C--D join opposite corners, any two on the same side cross, and three links cannot share two sides. Hint: ``Can you find a triangle in network 5? What does that say about its regions?''

\sub{Problem 6, page 6: the portal board}
As Grades 2--3 Problem 5.

\sub{Problems 7 and 8, page 7: six dots, then 10 and 11 dots}
Problem 7: 3 crossings (Figure 2). Nobody can do better: with 2 crossings, removing one link at each would leave at least 13 links on 6 dots untangled, but an untangled network on 6 dots has at most 12 (from regions: $3F\le 2E$ and $F=E-4$ give $E\le 12$; or because it can be continued as a game, which ends at 12). Only two links meet at a crossing, which is why one removal per crossing is enough. Follow-up: ``Could one link be in two crossings?'' (Removing it then clears both, which only helps.) Hint: ``What if you take away one link at each crossing?''

Problem 8: the second player wins with 10 dots (24 links) and the first with 11 (27). When nobody can move, the network is in one piece and every region is a triangle, or another link would fit; so $3\times\text{regions}=2\times\text{links}$, and with $\text{dots}-\text{links}+\text{regions}=2$ this gives $\text{links}=3\times\text{dots}-6$. Expect children to see 6, 9, 12 and ``three more links per dot'' first. That is a conjecture: putting a new dot inside a triangle of a finished game and joining it to the corners makes one finished game with three more links, not every game. Hint: ``When nobody can move, what do the regions look like?'' Follow-up: ``Why can't a finished game have a region with four dots round it?'' (If P, Q, R, S are in order round it, P--R and Q--S cannot both run outside the region without crossing, so one of them can still be drawn inside.)

\sub{Problem 9, page 8: six and seven dots on the portal board}
Both work at the printed positions (Figure 4): join each pair by its shortest route through the repeated board, and no two links cross. That is a property of these positions, not a general rule; for network 5 the shortest routes tie. Seven dots: 21 links, 10 through a portal, and the board fills with 14 triangles. Six dots: the same drawing without G, 15 links, 8 through a portal, 9 regions. Hint, only when needed: ``If the board repeated in every direction, as on Week 41's map, which copy of each dot is nearest?'' Extension for adults to discuss: eight dots all joined cannot fit even here, since $28>3\times8$.\par\needspace{8cm}
\begin{center}
\fig{portal_k6}{6 dots}\qquad\qquad\fig{portal_k7}{7 dots}\\[4pt]
{\small\textbf{Figure 4.} Six and seven dots at the printed positions, each pair joined by its shortest route; numbers as in Figure 3.}
\end{center}

\head{Sources, connections and checks}
\textbf{Mathematics.} Douglas B. West, \emph{Introduction to Graph Theory}, 2nd ed. (Prentice Hall, 2001), \S6.1: Euler's formula and $E\le3V-6$, $E\le2V-4$; section number from the outline, not consulted for this guide. K. Kuratowski, ``Sur le probl\`eme des courbes gauches en topologie'', \emph{Fundamenta Mathematicae} 15 (1930): the obstruction theorem, stated only; not consulted. The proofs here, the game analysis and the portal drawings are written for this guide and checked by computation.

\textbf{Puzzles.} Henry E. Dudeney, \emph{Amusements in Mathematics} (1917), ``Water, Gas, and Electricity'': problem 251, p.~73, solution pp.~200--201, as cited by the Wikipedia page ``Three utilities problem'' (consulted October 25, 2026); the book itself was not checked. That page reports that Dudeney's solution passes a pipe through a house, the loophole closed here by ``never lies across a dot''. ``Planarity'', Wikipedia (consulted October 25, 2026): John Tantalo's 2005 game, from a concept by Mary Radcliffe, in which players move dots one at a time until no links cross; the untangling format comes from it. Simon Tatham, ``Untangle'' (consulted October 25, 2026): the same puzzle, credited to Planarity.

\textbf{Library.} Week 56 proves $V-E+F=2$ on solids; here it becomes a tool for proving a drawing impossible. Week 41 built the portal board, whose seam marks these pages reuse. Week 40 studies crossings in knot drawings, and Week 10 counts the links at each island.

\textbf{Checks and what to record.} \texttt{check\_guide.py} recomputes every answer above from a rebuild of the final pages, including 300 random games on each printed game board, the bent launch link and the portal drawings of Figures 3 and 4 with their numbered exits; none of this is a physical rehearsal. Record the real kit time, whether clips hold four pipe-cleaner ends, whether games stop early, which networks children called stuck and why, the region rule in children's words, and which problems children want to continue.
